% from ega4-16-fr.tex
\oldpage[IV-4]{29}
finie, $\mathfrak{G}_{X/S}$ est un $\mathcal{O}_X$-Module quasi-cohérent; si
de plus $S$ est localement noethérien, $\mathfrak{G}_{X/S}$ est cohérent
(16.5.6).
\end{env}

\begin{env}[16.5.8]
\label{IV.16.5.8}
Supposons plus particulièrement que $\Omega_{X/S}^1$ soit un
$\mathcal{O}_X$-Module \emph{localement libre} (de rang fini) (ce qui sera le
cas lorsque $f$ est \emph{lisse} (17.2.3)); alors $\mathfrak{G}_{X/S}$ est un
$\mathcal{O}_X$-Module localement libre et a même rang que $\Omega_{X/S}^1$
en chaque point. De façon plus précise, supposons que $\Omega_{X/S}^1$ soit de
rang $n$ en un point $x$; il y a alors $n$ sections $s_i$ ($1\leq i\leq n$)
de $\mathcal{O}_X$ au-dessus d'un voisinage affine $U$ de $x$ telles que les
images canoniques des $ds_i$ dans
$\Omega_{X/S}^1\otimes_{\mathcal{O}_X}k(x)$ forment une base de ce
$k(x)$-espace vectoriel; en vertu du lemme de Nakayama, les germes
$(ds_i)_x=d(s_i)_x$ des $ds_i$ au point $x$ forment une base du
$\mathcal{O}_x$-module $(\Omega_{X/S}^1)_x$, donc, en restreignant $U$, on
peut supposer que les $ds_i$ forment une \emph{base} du
$\Gamma(U,\mathcal{O}_X)$-module $\Gamma(U,\Omega_{X/S}^1)$. Alors le
$\Gamma(U,\mathcal{O}_X)$-module $\Gamma(U,\mathfrak{G}_{X/S})$ est dual du
précédent; on note $(D_i)_{1\leq i\leq n}$ ou
$\left(\frac{\partial}{\partial s_i}\right)_{1\leq i\leq n}$ la base duale de
$(ds_i)_{1\leq i\leq n}$, de sorte que, par (16.5.3), on a
\[
  D_is_j=\langle D_i,ds_j\rangle
  =\left\langle\frac{\partial}{\partial s_i},ds_j\right\rangle
  =\delta_{ij}\qquad\text{(indice de Kronecker).}
  \tag{16.5.8.1}
  \label{IV.16.5.8.1}
\]
Toute $\Gamma(S,\mathcal{O}_S)$-dérivation de la
$\Gamma(S,\mathcal{O}_S)$-algèbre $\Gamma(U,\mathcal{O}_X)$ s'écrit donc
d'une seule manière
\[
  D=\sum_{i=1}^n a_iD_i=\sum_{i=1}^n a_i\frac{\partial}{\partial s_i}
\]
où les $a_i$ ($1\leq i\leq n$) sont des sections de $\mathcal{O}_X$ au-dessus
de $U$. Pour toute section $g\in\Gamma(U,\mathcal{O}_X)$, si l'on pose
$dg=\sum_{i=1}^n c_i ds_i$, on a
$c_i=\langle D_i,dg\rangle=D_ig$ en vertu de (16.5.8.1), autrement dit
\[
  dg=\sum_{i=1}^n(D_ig)ds_i
  =\sum_{i=1}^n\frac{\partial g}{\partial s_i}ds_i.
  \tag{16.5.8.2}
  \label{IV.16.5.8.2}
\]
\end{env}

\begin{env}[16.5.9]
\label{IV.16.5.9}
Soient $D_1$, $D_2$ deux $S$-dérivations de $\mathcal{O}_X$. Pour tout ouvert
$U$ de $X$, si $D_1^U$, $D_2^U$ sont les dérivations correspondantes de
l'anneau $\Gamma(U,\mathcal{O}_X)$, le \emph{crochet}
\[
  [D_1^U,D_2^U]=D_1^U\circ D_2^U-D_2^U\circ D_1^U
\]
est aussi une dérivation de cet anneau, donc le
$\psi^*(\mathcal{O}_S)$-endomorphisme de $\mathcal{O}_X$
\[
  [D_1,D_2]=D_1\circ D_2-D_2\circ D_1
  \tag{16.5.9.1}
  \label{IV.16.5.9.1}
\]
est encore une $S$-dérivation; comme on vérifie aussitôt que ce crochet
vérifie l'identité de Jacobi, on voit qu'on a ainsi défini sur
$\operatorname{D\acute{e}r}_S(\mathcal{O}_X,\mathcal{O}_X)$ une structure de
$\Gamma(S,\mathcal{O}_S)$-algèbre de Lie. Comme la définition de cette
structure commute à la restriction à un ouvert de $X$, on voit que
$\mathfrak{G}_{X/S}$ est canoniquement muni d'une structure de
$\psi^*(\mathcal{O}_S)$-Algèbre de Lie. On notera que l'application
$(D_1,D_2)\mapsto[D_1,D_2]$ n'est pas
$\Gamma(X,\mathcal{O}_X)$-bilinéaire.
\end{env}

\begin{env}[16.5.10]
\label{IV.16.5.10}
Pour tout changement de base $g:S'\to S$, si l'on pose $X'=X\times_S S'$, on
a vu (16.4.5) que l'on a un isomorphisme canonique
\[
  \Omega_{X/S}^1\otimes_S S'\xrightarrow{\sim}\Omega_{X'/S'}^1
  \tag{16.5.10.1}
  \label{IV.16.5.10.1}
\]
